% ═══════════════════════════════════════════════
%  MỞ ĐẦU
% ═══════════════════════════════════════════════

\section{MỞ ĐẦU}
\label{ch:intro}

% ───────────────────────────────────────────────
\subsection{Lý do chọn đề tài}
% ───────────────────────────────────────────────

\indent Trò chơi điện tử giữ một vị trí khác thường trong nghiên cứu trí tuệ nhân tạo:
chúng đồng thời là phòng thí nghiệm và là sản phẩm. Một môi trường game cho phép chạy
hàng triệu lần thử với chi phí gần bằng không và không có rủi ro nào, điều mà robot
thật hay hệ thống giao thông thật không bao giờ cho phép. Ở chiều ngược lại, chính
ngành công nghiệp game có nhu cầu thương mại trực tiếp với những nhân vật máy biết cư
xử thuyết phục. Hiếm lĩnh vực nào mà quãng đường từ một kết quả nghiên cứu tới một
tính năng bán được lại ngắn đến thế.

\indent Học tăng cường là hướng tiếp cận ăn khớp nhất với bối cảnh này. Nó không đòi
hỏi tập dữ liệu gán nhãn sẵn, thứ gần như không thể có cho câu hỏi ``nước đi nào là
đúng trong tình huống này'', mà để tác nhân tự rút ra chiến lược từ hậu quả của chính
những hành động nó đã thử.

\indent Khả năng của học tăng cường đã được chứng minh qua hàng loạt thành tựu mang
tính bước ngoặt. Một hệ thống học tăng cường có thể chơi thành thạo hàng chục trò chơi
điện tử cổ điển chỉ từ hình ảnh thô của màn hình \cite{mnih2015humanlevel}, có thể
đánh bại kỳ thủ cờ vây hàng đầu thế giới \cite{silver2016mastering}, hay phối hợp
nhiều tác nhân để chiến thắng trong những trò chơi chiến thuật phức tạp
\cite{openai2019dota, vinyals2019grandmaster}. Điểm chung của các hệ thống này là
không một hành vi nào trong số đó được lập trình sẵn: chúng đều nảy sinh từ quá trình
thử và sai của chính tác nhân.

\indent Tuy nhiên, ấn tượng mà những thành tựu ấy để lại dễ khiến người ta hiểu sai về
nơi khó khăn thực sự nằm ở đâu. Các thuật toán học tăng cường hiện đã được cài đặt sẵn,
kiểm chứng kỹ và đóng gói trong những thư viện chín muồi; gọi ra một trình huấn luyện
PPO chỉ tốn một dòng lệnh. Thứ không có sẵn, và cũng là thứ chiếm phần lớn công sức của
một dự án học tăng cường thực tế, chính là \textit{môi trường huấn luyện}. Người xây
dựng phải tự trả lời: tác nhân nhìn thấy gì, hành động ra sao, được thưởng phạt theo
tiêu chí nào, tình huống được đặt lại thế nào sau mỗi lượt, và làm sao để toàn bộ hệ
thống chạy đủ nhanh mà thu về hàng triệu bước kinh nghiệm. Mỗi câu hỏi đó là một quyết
định thiết kế, và một quyết định sai thường không báo lỗi: chương trình vẫn chạy, đường
cong vẫn vẽ ra, chỉ là tác nhân không bao giờ học được điều ta muốn.

\indent Các môi trường hai chiều đơn giản hoặc các bộ dữ liệu chuẩn thường dùng trong
nghiên cứu học thuật tuy thuận tiện nhưng lại giấu đi đúng phần khó nói trên: chúng đã
được đặc tả sẵn từ trước, người dùng chỉ việc nạp vào. Một môi trường ba chiều với mô
phỏng vật lý chân thực, quan sát một phần qua cảm biến và tình huống thay đổi liên tục
mới bộc lộ đầy đủ các bài toán thiết kế, đồng thời gần với bài toán phát triển game
thương mại hơn. Việc thiếu vắng những môi trường như vậy ở dạng dùng lại được - có tài
liệu, có công cụ, có cấu hình huấn luyện đi kèm - là một khoảng trống đáng chú ý, và
cũng là khoảng trống mà đồ án này nhắm tới.

\indent Nền tảng Unity kết hợp với bộ công cụ ML-Agents cung cấp một hệ sinh thái đủ
mạnh để giải quyết vấn đề trên. Unity là một trong những công cụ phát triển game phổ
biến nhất hiện nay, hỗ trợ mô phỏng vật lý và dựng hình ba chiều ở mức chuyên nghiệp;
còn ML-Agents nối thẳng một cảnh Unity đang chạy với các thuật toán học tăng cường
được cài đặt trong hệ sinh thái Python.
Sự kết hợp này biến Unity thành một nền tảng lý tưởng để thiết kế các môi trường huấn
luyện đa dạng, từ bài toán đối kháng, hợp tác cho đến điều hướng tránh vật cản trong
không gian ba chiều.

\indent Xuất phát từ những lý do trên, tôi lựa chọn đề tài ``Học tăng cường
để tối ưu hóa khả năng tương tác của AI trong game'' với trọng tâm là \textbf{nghiên
cứu và xây dựng môi trường huấn luyện} cho tác nhân học tăng cường trên nền tảng Unity
ML-Agents: ba môi trường ba chiều thuộc ba dạng bài toán khác nhau, kèm toàn bộ quá
trình thiết kế quan sát - hành động - phần thưởng, các công cụ hỗ trợ và cách đóng gói
để người khác dùng lại được.

\indent Việc huấn luyện năm thuật toán PPO, SAC, MA-POCA, MAPPO và DQN trên cả ba môi
trường vẫn được tiến hành đầy đủ, nhưng vai trò của nó trong đề tài là \textit{kiểm
chứng môi trường} chứ không phải xếp hạng thuật toán. Lý do là một môi trường huấn
luyện không thể được nghiệm thu chỉ bằng cách nhìn vào mã nguồn: cách duy nhất biết
đặc tả của nó đúng hay sai là thả nhiều thuật toán khác nhau vào và quan sát chúng cư
xử ra sao. Một môi trường mà mọi thuật toán đều thất bại thì nhiều khả năng lỗi nằm ở
phần thưởng chứ không ở thuật toán; ngược lại một môi trường mà thuật toán nào cũng đạt
điểm tối đa thì nhiều khả năng đang tồn tại lối tắt. Chỉ khi các thuật toán tách ra
thành thứ hạng có thể giải thích được thì môi trường mới thực sự đo được thứ nó định
đo. Đích đến của đề tài, vì vậy, là để lại một bộ môi trường đã được kiểm chứng và đóng
gói cho các nghiên cứu và dự án phát triển game về sau.

% ───────────────────────────────────────────────
\subsection{Mục tiêu nghiên cứu}

\indent Mục tiêu tổng quát của đề tài là nghiên cứu quy trình thiết kế, xây dựng và
kiểm chứng môi trường huấn luyện cho tác nhân học tăng cường trong game ba chiều trên
nền tảng Unity ML-Agents, đồng thời tạo ra một bộ môi trường dùng lại được. Từ mục tiêu
tổng quát đó, đề tài hướng đến năm mục tiêu cụ thể, trong đó ba mục tiêu đầu là trọng
tâm.

\begin{itemize}
    \item \textbf{Xây dựng bộ môi trường huấn luyện.} Đề tài xây dựng ba môi trường
    game ba chiều có tính chất khác biệt để bao phủ ba dạng bài toán học tăng cường
    phổ biến: Football Table mô phỏng một bàn bi lắc ba chiều, đại diện cho bài toán đối
    kháng một đối một với không gian hành động liên tục; Capture The Flag đòi hỏi nhiều
    tác nhân phối hợp với nhau để cùng hoàn thành một mục tiêu chung; và Cross The Road
    buộc tác nhân điều hướng vượt qua dòng phương tiện di chuyển liên tục để đến đích.
    Ba dạng bài toán này được chọn vì chúng đặt ra ba nhóm bài toán thiết kế khác hẳn
    nhau: gán công lao trong tự chơi, phần thưởng nhóm trong hợp tác, và phần thưởng
    thưa trong điều hướng.

    \item \textbf{Nghiên cứu quy trình thiết kế quan sát, hành động và phần thưởng.} Đây
    là phần lõi về mặt học thuật của đề tài. Với mỗi môi trường, đề tài trình bày không
    chỉ phương án cuối cùng mà cả quá trình đi đến nó: những đặc tả đã thử và thất bại,
    dấu hiệu nhận biết một hàm phần thưởng bị hỏng, và cách sửa. Mục tiêu là rút ra được
    những nguyên tắc thiết kế dùng lại được, chứ không chỉ ba cấu hình chạy được.

    \item \textbf{Đóng gói môi trường thành sản phẩm dùng lại được.} Một môi trường chỉ
    có giá trị với người khác khi họ cài được nó vào dự án của mình. Đề tài đóng gói cả
    ba môi trường thành một gói Unity Package Manager kèm cấu hình huấn luyện, bộ công cụ
    Editor (nhân khu vực huấn luyện, build headless, sinh biến thể không gian hành động)
    và một trình hướng dẫn cài đặt tự kiểm tra trạng thái từng bước.

    \item \textbf{Kiểm chứng môi trường bằng thực nghiệm đa thuật toán.} Đề tài huấn
    luyện năm thuật toán PPO, SAC, MA-POCA, MAPPO và DQN trên từng môi trường trong cùng
    điều kiện phần cứng và cùng ngân sách bước, tạo thành một lưới đầy đủ mười lăm lần
    chạy. Mục đích là xác nhận ba tính chất của một môi trường đạt yêu cầu - tác nhân học
    được, môi trường phân biệt được các hướng tiếp cận, và không tồn tại lối tắt thoái
    hóa - đồng thời đo đặc tính của từng môi trường theo bốn tiêu chí định lượng: mức
    năng lực đạt được, tốc độ hội tụ, độ ổn định huấn luyện và độ dài lượt chơi.

    \item \textbf{Xây dựng ứng dụng trình diễn để quan sát hành vi.} Biểu đồ phần thưởng
    không cho biết tác nhân đang làm gì, trong khi nhiều lỗi đặc tả chỉ lộ ra khi nhìn
    hành vi thật. Đề tài vì vậy xây dựng một ứng dụng cho phép nạp trực tiếp các mô hình
    đã huấn luyện dưới định dạng ONNX, tự điều khiển hoặc chơi cùng tác nhân để đối chiếu
    độ khó, và tra cứu ngay số liệu huấn luyện trên cùng giao diện.
\end{itemize}

% ───────────────────────────────────────────────
\subsection{Đối tượng và phạm vi nghiên cứu}

\indent Đối tượng nghiên cứu chính của đề tài là \textbf{môi trường huấn luyện học tăng
cường trong game ba chiều}: cấu trúc của nó, các quyết định thiết kế phải đưa ra khi xây
dựng, và cách kiểm chứng xem nó đã được đặc tả đúng hay chưa. Cụ thể hóa thành ba môi
trường được xây dựng trên nền tảng Unity: môi trường đối kháng Football Table, môi
trường hợp tác Capture The Flag và môi trường điều hướng Cross The Road.

\indent Đối tượng thứ hai là bộ công cụ phục vụ việc xây dựng và vận hành môi trường:
nền tảng Unity ML-Agents, cơ chế trainer mở rộng, giao diện lập trình cấp thấp bằng
Python, cùng các công cụ do đề tài tự viết để nhân khu vực huấn luyện, dựng bản build
không đồ hoạ và đóng gói môi trường thành gói phân phối.

\indent Năm thuật toán PPO (tối ưu chính sách vùng lân cận), SAC (actor-critic cực đại
hóa entropy), MA-POCA (gán công lao cho đa tác nhân), MAPPO (mở rộng PPO cho đa tác nhân
với chia sẻ tham số) và DQN (dựa trên giá trị với bộ đệm phát lại) tham gia đề tài với
tư cách \textit{công cụ đo}: chúng là những phép thử độc lập dùng để dò đặc tính của môi
trường. Đề tài không đặt mục tiêu cải tiến hay đề xuất thuật toán mới.

\indent Về phạm vi nội dung, đề tài tập trung vào thiết kế môi trường ba chiều, xây
dựng hệ thống quan sát, hành động, phần thưởng, quy trình huấn luyện và cách đóng gói
môi trường để dùng lại. Đề tài không đi sâu vào các khía cạnh tối ưu hóa đồ họa, âm
thanh hay hoàn thiện sản phẩm game thương mại, cũng không đặt mục tiêu đạt hiệu năng
cao nhất có thể trên từng môi trường - ngân sách bước được chọn ở mức đủ để các đường
cong bộc lộ đặc tính của bài toán, không phải để ép mỗi thuật toán tới giới hạn của nó.
Ba môi trường được lựa chọn nhằm bao phủ ba dạng bài toán học tăng cường điển hình là
đối kháng một đối một, hợp tác đa tác nhân và điều hướng đơn tác nhân.

\indent Về phạm vi thuật toán, năm thuật toán PPO, SAC, MA-POCA, MAPPO và DQN được chọn
dựa trên mức độ phổ biến trong nghiên cứu, tính sẵn có trong hệ sinh thái ML-Agents và
sự đa dạng về mặt kiến trúc, khi chúng đại diện cho năm hướng tiếp cận khác nhau lần
lượt là on-policy đơn tác nhân, actor-critic off-policy, đa tác nhân dựa trên attention,
đa tác nhân dựa trên chia sẻ tham số và phương pháp dựa trên giá trị. Vì DQN đòi hỏi
không gian hành động rời rạc, môi trường Football Table vốn dùng hành động liên tục
phải được dựng lại thành một bản riêng có hành động rời rạc dành cho lần chạy này.
Toàn bộ thực nghiệm được tiến hành trên một cấu hình máy
tính cố định để bảo đảm tính công bằng khi so sánh, chi tiết cấu hình này được trình
bày trong Mục~\ref{sub:hardware}.

% ───────────────────────────────────────────────
\subsection{Tổng quan tài liệu và khoảng trống nghiên cứu}

\indent Nghiên cứu về học tăng cường trong game đã trải qua nhiều giai đoạn phát triển.
Mnih và cộng sự \cite{mnih2015humanlevel} đề xuất DQN và lần đầu chứng minh rằng học
tăng cường sâu có thể đạt trình độ con người trong các trò chơi điện tử cổ điển chỉ từ
dữ liệu điểm ảnh thô, mở ra hướng ứng dụng học tăng cường trực tiếp vào môi trường
game; đây cũng là lý do DQN được giữ lại trong đề tài như một mốc so sánh lịch sử cho
nhóm phương pháp dựa trên giá trị. Tiếp nối
thành tựu đó, Silver và cộng sự \cite{silver2016mastering, silver2017mastering} đạt
bước đột phá với các hệ thống chơi cờ vây ở trình độ kiện tướng thế giới, khi kết hợp
học tăng cường với mạng nơ-ron sâu và tìm kiếm cây.

\indent Ở quy mô lớn hơn với môi trường phức tạp và nhiều tác nhân, các hệ thống chơi
trò chơi chiến thuật thời gian thực \cite{openai2019dota, vinyals2019grandmaster} đã
cho thấy học tăng cường có thể xử lý những bài toán với không gian trạng thái khổng
lồ, nhiều tác nhân hoạt động đồng thời và đòi hỏi chiến lược dài hạn. Về phía thuật
toán, Schulman và cộng sự \cite{schulman2017proximal} đề xuất PPO, một phương pháp
on-policy chấp nhận đánh đổi hiệu quả mẫu để lấy sự ổn định và tính dễ triển khai,
trong khi Haarnoja và cộng sự
\cite{haarnoja2018soft} giới thiệu SAC với cơ chế điều hòa entropy giúp cân bằng giữa
khám phá và khai thác, đặc biệt hiệu quả với không gian hành động liên tục. Trong lĩnh
vực đa tác nhân, Cohen và cộng sự \cite{cohen2022mapoca} phát triển MA-POCA với cơ chế
phê bình tập trung nhằm giải quyết bài toán gán công lao trong môi trường hợp tác,
trong khi Yu và cộng sự \cite{yu2022surprising} chứng minh rằng MAPPO - một mở rộng đơn
giản của PPO sang nhiều tác nhân với chia sẻ tham số - đạt hiệu năng cạnh tranh hoặc
vượt trội trên nhiều bộ benchmark hợp tác chuẩn.

\indent Về mặt công cụ, Juliani và cộng sự \cite{juliani2020unity} giới thiệu Unity
ML-Agents như một hệ sinh thái mã nguồn mở tích hợp sẵn khả năng xây dựng môi trường
ba chiều với mô phỏng vật lý chân thực, hỗ trợ huấn luyện tác nhân trực tiếp và được
sử dụng rộng rãi trong cả nghiên cứu lẫn phát triển game.

\indent Điểm chung của toàn bộ các công trình trên là chúng \textit{tiêu thụ} môi
trường chứ không bàn về việc tạo ra môi trường. Bộ benchmark đã được ai đó đặc tả sẵn -
lưới hai chiều, vector trạng thái rút gọn, động lực học đơn giản - và đóng góp khoa học
nằm ở phía thuật toán. Chính vì vậy, tài liệu về việc \textit{thiết kế} một môi trường
học tăng cường lại mỏng một cách bất ngờ so với tài liệu về thuật toán: những câu hỏi
như nên đưa đại lượng nào vào quan sát, phần thưởng dày hay thưa, phạt thời gian bao
nhiêu là đủ, hay khi nào một hàm phần thưởng đang vô tình trừng phạt chính hành vi ta
muốn dạy, hầu như chỉ được truyền lại dưới dạng kinh nghiệm rời rạc trong các diễn đàn
kỹ thuật.

\indent Khoảng trống mà đề tài nhắm tới nằm đúng ở đó, và có hai mặt. Mặt thứ nhất là
\textit{tri thức thiết kế}: quá trình đi từ một ý tưởng trò chơi tới một đặc tả mà tác
nhân học được rất hiếm khi được ghi lại đầy đủ, đặc biệt là các phương án đã thất bại -
vốn là phần dạy được nhiều nhất. Mặt thứ hai là \textit{hiện vật dùng lại được}: các
môi trường ba chiều đi kèm tài liệu, cấu hình huấn luyện và công cụ cài đặt ở dạng một
nhóm phát triển khác có thể lấy về dùng ngay gần như không có, nên mỗi dự án lại phải
dựng lại từ đầu. Đề tài này lấp cả hai mặt: ghi lại toàn bộ quá trình thiết kế ba môi
trường thuộc ba dạng bài toán khác nhau, rồi đóng gói chúng thành sản phẩm phân phối
được. Việc đặt PPO, SAC, MA-POCA, MAPPO và DQN lên cùng một bàn cân - điều gần như
không thực hiện được chỉ bằng cách đối chiếu tài liệu đã công bố - là phương tiện để
kiểm chứng ba môi trường đó, và cũng là một kết quả phụ có ích cho người làm game.

% ───────────────────────────────────────────────
\subsection{Phương pháp nghiên cứu}
\indent Đề tài sử dụng kết hợp ba nhóm phương pháp nghiên cứu bổ trợ lẫn nhau. Phương
pháp nghiên cứu lý thuyết được dùng để tổng hợp và phân tích các tài liệu khoa học về
học tăng cường, học tăng cường sâu và năm thuật toán PPO, SAC, MA-POCA, MAPPO và DQN,
dựa trên các bài báo gốc cùng tài liệu kỹ thuật của Unity ML-Agents; đây là cơ sở cho
việc thiết kế môi trường và lựa chọn tham số huấn luyện.

\indent Phương pháp thực nghiệm được sử dụng để xây dựng ba môi trường game ba chiều
trong Unity với đầy đủ hệ thống vật lý, quan sát, hành động và phần thưởng, sau đó
triển khai và huấn luyện từng thuật toán trên mỗi môi trường trong cùng điều kiện phần
cứng và số bước huấn luyện. Trong quá trình này, dữ liệu huấn luyện được thu thập và
trực quan hóa theo thời gian thực bằng công cụ TensorBoard.

\indent Cuối cùng, phương pháp so sánh và đánh giá được dùng để đối chiếu hiệu năng
của các thuật toán dựa trên bốn tiêu chí định lượng là mức năng lực đạt được, tốc độ
hội tụ, độ ổn định huấn luyện và độ dài lượt chơi, kết hợp với phân tích
định tính về mức độ phù hợp của từng thuật toán với từng dạng môi trường.

% ───────────────────────────────────────────────
\subsection{Ý nghĩa khoa học và thực tiễn}

\indent Về mặt khoa học, đề tài đề xuất và hiện thực hóa một quy trình xây dựng môi
trường game ba chiều phục vụ huấn luyện học tăng cường trên nền tảng Unity, kèm theo
một cách nghiệm thu môi trường bằng thực nghiệm đa thuật toán dựa trên ba tiêu chí học
được - phân biệt được - không có lối tắt. Quy trình này dùng lại được cho bất kỳ môi
trường nào khác, không riêng ba môi trường của đề tài. Phần ghi chép về các phương án
đặc tả đã thất bại cũng có giá trị riêng, vì đây là loại tri thức thường bị bỏ lại
ngoài các công bố khoa học dù nó chiếm phần lớn thời gian thực hiện. Ngoài ra, quy
trình bổ sung thuật toán mới vào ML-Agents dưới dạng trainer mở rộng, cùng cách kết nối
một thuật toán off-policy với Unity thông qua giao diện lập trình cấp thấp bằng Python,
cũng được trình bày chi tiết, mở ra một hướng đi có thể tái sử dụng cho các nghiên cứu
cần thuật toán nằm ngoài bộ tích hợp sẵn của Unity.

\indent Về mặt thực tiễn, ba môi trường ba chiều được xây dựng trong đề tài được đóng
gói ở dạng cài đặt được ngay vào một dự án Unity bất kỳ, nên có thể dùng trực tiếp làm
nền tảng thử nghiệm trí tuệ nhân tạo mà không phải dựng lại từ đầu - vốn là phần tốn
kém nhất của một dự án học tăng cường. Quy trình thiết kế hàm phần thưởng và không gian
quan sát trong môi trường ba chiều được trình bày trong đề tài có giá trị tham khảo cho
các nhóm phát triển game có nhu cầu tích hợp trí tuệ nhân tạo. Kết quả so sánh năm thuật
toán, tuy đóng vai trò kiểm chứng trong đề tài, cũng giúp các nhà phát triển có cơ sở
lựa chọn thuật toán phù hợp với từng cơ chế game mà không phải tiến hành lại toàn bộ
thực nghiệm.

% ───────────────────────────────────────────────
\subsection{Bố cục đồ án}
\indent Ngoài phần Mở đầu và Kết luận, nội dung chính của đồ án được tổ chức thành bốn
chương. Chương~\ref{ch:theory} trình bày cơ sở lý thuyết, bao gồm tổng quan về trí tuệ
nhân tạo trong game, nền tảng lý thuyết của học tăng cường với quá trình quyết định
Markov, chính sách và hàm giá trị, các kiến trúc học tăng cường sâu, cùng đặc điểm của
năm thuật toán PPO, SAC, MA-POCA, MAPPO và DQN; chương này cũng giới thiệu nền tảng Unity
ML-Agents và các giới hạn tích hợp liên quan đến việc triển khai thuật toán tùy chỉnh.

\indent Chương~\ref{ch:design} trình bày chi tiết quá trình thiết kế ba môi trường
game ba chiều trong Unity, bao gồm không gian quan sát, không gian hành động và hàm
phần thưởng của từng môi trường; đây là chương trọng tâm của đồ án. Chương này cũng mô
tả phương pháp đưa các thuật toán nằm ngoài bộ huấn luyện tích hợp sẵn của ML-Agents
vào thực nghiệm, thiết lập thực nghiệm và các tiêu chí đánh giá được sử dụng.
Chương~\ref{ch:results} dùng mười lăm lần huấn luyện để nghiệm thu ba môi trường đó,
trình bày kết quả của từng thuật toán trên từng môi trường qua biểu đồ và bảng số liệu,
đối chiếu chúng để xác nhận môi trường học được, phân biệt được và không có lối tắt,
đồng thời rút ra nhận xét về mức độ phù hợp giữa thuật toán và dạng bài toán.

\indent Cuối cùng, Chương~\ref{ch:app} trình bày quá trình xây dựng ứng dụng trình
diễn tương tác: kiến trúc hệ thống lựa chọn chế độ chơi, giao diện menu, cơ chế nạp
mô hình ONNX tại thời gian chạy, đường ống đưa số liệu huấn luyện từ TensorBoard vào
giao diện trong game, cách ánh xạ ba chế độ chơi (người chơi, tác nhân tự chơi, người
chơi cùng tác nhân) vào từng môi trường, cùng các
vấn đề kỹ thuật phát sinh và kết quả kiểm thử ứng dụng.
